\hypertarget{deep-object-model-virtualization}{%
\chapter{DEEP OBJECT MODEL \&
VIRTUALIZATION}\label{deep-object-model-virtualization}}

Understanding the ``C++ Object Model'' distinguishes a user from a
master. This section explains what the compiler generates for your
classes.

\hypertarget{the-cost-of-polymorphism-vptr-vtable}{%
\subsection{2.5.1 The Cost of Polymorphism (vptr \&
vtable)}\label{the-cost-of-polymorphism-vptr-vtable}}

Every class with \emph{at least one} virtual function has a hidden
overhead.

\begin{enumerate}
\def\labelenumi{\arabic{enumi}.}
\tightlist
\item
  \textbf{vptr (Virtual Pointer)}: A hidden member added to the
  \emph{layout} of the object.
\item
  \textbf{vtable (Virtual Table)}: A static table of function pointers
  for that class.
\end{enumerate}

\begin{lstlisting}[language={C++}]
class Base {
    int data;
    virtual void func() {}
};
// sizeof(Base) = sizeof(int) + sizeof(void*) + padding
// On 64-bit: 4 bytes (int) + 4 bytes (padding) + 8 bytes (ptr) = 16 bytes
\end{lstlisting}

\hypertarget{multiple-inheritance-thunks}{%
\subsection{2.5.2 Multiple Inheritance \&
Thunks}\label{multiple-inheritance-thunks}}

When inheriting from multiple classes, pointer arithmetic gets tricky.

\begin{lstlisting}[language={C++}]
class A { int a; virtual void f() {} };
class B { int b; virtual void g() {} };
class C : public A, public B { int c; };

C obj;
A* pa = &obj; // Points to start of obj
B* pb = &obj; // Points to obj + sizeof(A) !!
\end{lstlisting}

\begin{itemize}
\tightlist
\item
  \textbf{Thunk}: A small piece of assembly code generated by the
  compiler to adjust the \passthrough{\lstinline!this!} pointer when
  calling a virtual function from a base class pointer that isn't at
  offset 0.
\end{itemize}

\hypertarget{virtual-inheritance-the-diamond-problem}{%
\subsection{2.5.3 Virtual Inheritance (The Diamond
Problem)}\label{virtual-inheritance-the-diamond-problem}}

\begin{lstlisting}[language={C++}]
class Top { int t; };
class Left : virtual public Top { int l; };
class Right : virtual public Top { int r; };
class Bottom : public Left, public Right { int b; };
\end{lstlisting}

To solve the Diamond Problem (Top appearing twice),
\passthrough{\lstinline!virtual!} inheritance ensures
\passthrough{\lstinline!Top!} is shared. * \textbf{Cost}:
\passthrough{\lstinline!Left!} and \passthrough{\lstinline!Right!} now
contain a \textbf{vbptr} (Virtual Base Pointer) pointing to the shared
\passthrough{\lstinline!Top!} instance. Accessing members of
\passthrough{\lstinline!Top!} becomes slower (indirection).

\hypertarget{alignment-padding-rules}{%
\subsection{2.5.4 Alignment \& Padding
Rules}\label{alignment-padding-rules}}

CPU reads are efficient at specific addresses (multiples of 4 or 8).
Compilers insert ``padding bytes''.

\textbf{Rule}: A member of size \(N\) must sit at an offset divisible by
\(N\).

\begin{lstlisting}[language={C++}]
struct Mixed {
    char a;     // 1 byte
                // 3 bytes PADDING
    int b;      // 4 bytes
    short c;    // 2 bytes
                // 6 bytes PADDING (to align structure size to 8)
};
// sizeof(Mixed) = 16 (on 64-bit)
\end{lstlisting}

\textbf{Optimization}: Sort members by size (Largest to Smallest) to
minimize padding.

\begin{center}\rule{0.5\linewidth}{0.5pt}\end{center}
